% sec:dg-curvature — Curvature Tensors
\section{Curvature Tensors}
\label{sec:dg-curvature}

Parallel-transport a vector around a closed loop on a sphere, and it comes
back rotated.  On a flat surface, this never happens --- the vector returns to
its starting orientation regardless of the path.  The discrepancy between the
initial and final vectors encodes the curvature enclosed by the loop, and the
Riemann tensor is the local, infinitesimal version of this statement: it
measures how much a vector rotates when parallel-transported around a
vanishingly small loop.  \Cref{sec:dg-geodesic} promised a measure of how the
gravitational field curves nearby geodesics apart --- the Riemann tensor is
that measure.
\physnote{Geodesic deviation --- the relative acceleration of neighbouring
freely falling particles --- is proportional to the Riemann tensor.  This is
how tidal forces arise in general relativity: not as forces, but as curvature.}

\begin{figure}[htbp]
  \centering
  % parallel-transport-loop.tex — Parallel transport around a spherical triangle
% Inputable TikZ fragment for fig:parallel-transport-loop
% Requires: tikz with calc, arrows.meta (loaded by ehbook.cls)
%
% Projection: orthographic, 30° elevation tilt (matches fig:sphere-geodesics).
%
% Triangle vertices (rotated to face the viewer):
%   A = (θ,φ) = (π/2, 225°)   equator, lower-left
%   B = (0, 0)                  north pole
%   C = (π/2, 315°)            equator, lower-right
% The triangle subtends solid angle π/2 (one-eighth of the sphere).
%
% A vector initially pointing north at A (upward on screen) is parallel-
% transported around A→B→C→A.  It returns pointing along the equator
% (rightward on screen) — rotated by π/2.
%
\begin{tikzpicture}[
    >={Stealth[length=5pt]},
    font=\footnotesize,
  ]

  % ── Parameters ──────────────────────────────────────────────────────
  \def\R{1.8}
  \pgfmathsetmacro{\tilt}{30}
  \pgfmathsetmacro{\costilt}{cos(\tilt)}
  \pgfmathsetmacro{\sintilt}{sin(\tilt)}
  \pgfmathsetmacro{\eqb}{\R*\sintilt}     % equator semi-minor axis

  % ── Projected vertex positions ─────────────────────────────────────
  % A = (90°, 225°): embed (−.707, −.707, 0)
  \pgfmathsetmacro{\Ax}{-0.7071*\R}
  \pgfmathsetmacro{\Ay}{-0.7071*\R*\sintilt}
  % B = (0°, 0°): embed (0, 0, 1)
  \pgfmathsetmacro{\Bx}{0}
  \pgfmathsetmacro{\By}{\R*\costilt}
  % C = (90°, 315°): embed (.707, −.707, 0)
  \pgfmathsetmacro{\Cx}{0.7071*\R}
  \pgfmathsetmacro{\Cy}{-0.7071*\R*\sintilt}

  % ── Sphere outline and shading ──────────────────────────────────────
  \shade[ball color=EHMarginGray!8, draw=EHMarginGray, thin]
    (0,0) circle (\R);

  % ── Equator (faint, near side only: φ = 180° to 360°) ─────────────
  \draw[EHMarginGray, thin, opacity=0.15]
    ({-\R}, 0) arc (180:360:{\R} and {\eqb});

  % ── Triangle fill (faint gold) ─────────────────────────────────────
  % Fill uses sampled meridian curves + equator arc
  \fill[EHAccretionGold, opacity=0.08]
    % A→B along meridian φ=225°
    ({-0.7071*\R}, {-0.7071*\R*\sintilt})
    ({-0.7071*\R*sin(75)}, {(-0.7071*sin(75)*\sintilt + cos(75)*\costilt)*\R})
    ({-0.7071*\R*sin(60)}, {(-0.7071*sin(60)*\sintilt + cos(60)*\costilt)*\R})
    ({-0.7071*\R*sin(45)}, {(-0.7071*sin(45)*\sintilt + cos(45)*\costilt)*\R})
    ({-0.7071*\R*sin(30)}, {(-0.7071*sin(30)*\sintilt + cos(30)*\costilt)*\R})
    ({-0.7071*\R*sin(15)}, {(-0.7071*sin(15)*\sintilt + cos(15)*\costilt)*\R})
    (0, {\R*\costilt})
    % B→C along meridian φ=315°
    ({0.7071*\R*sin(15)}, {(-0.7071*sin(15)*\sintilt + cos(15)*\costilt)*\R})
    ({0.7071*\R*sin(30)}, {(-0.7071*sin(30)*\sintilt + cos(30)*\costilt)*\R})
    ({0.7071*\R*sin(45)}, {(-0.7071*sin(45)*\sintilt + cos(45)*\costilt)*\R})
    ({0.7071*\R*sin(60)}, {(-0.7071*sin(60)*\sintilt + cos(60)*\costilt)*\R})
    ({0.7071*\R*sin(75)}, {(-0.7071*sin(75)*\sintilt + cos(75)*\costilt)*\R})
    ({0.7071*\R}, {-0.7071*\R*\sintilt})
    % C→A along equator (arc approximated by returning)
    -- ({0.7071*\R}, {-0.7071*\R*\sintilt})
       arc (315:225:{\R} and {\eqb})
    -- cycle;

  % ── Triangle edge: A→B (meridian φ=225°, θ from 90° to 0°) ────────
  % x = −0.707·R·sinθ,  y = R·(−0.707·sinθ·sintilt + cosθ·costilt)
  \draw[EHAccretionGold, thick] plot[smooth, tension=0.55] coordinates {
    ({-0.7071*\R},               {(-0.7071*\sintilt + 0*\costilt)*\R})
    ({-0.7071*\R*sin(75)}, {(-0.7071*sin(75)*\sintilt + cos(75)*\costilt)*\R})
    ({-0.7071*\R*sin(60)}, {(-0.7071*sin(60)*\sintilt + cos(60)*\costilt)*\R})
    ({-0.7071*\R*sin(45)}, {(-0.7071*sin(45)*\sintilt + cos(45)*\costilt)*\R})
    ({-0.7071*\R*sin(30)}, {(-0.7071*sin(30)*\sintilt + cos(30)*\costilt)*\R})
    ({-0.7071*\R*sin(15)}, {(-0.7071*sin(15)*\sintilt + cos(15)*\costilt)*\R})
    (0,                          {\R*\costilt})
  };

  % ── Triangle edge: B→C (meridian φ=315°, θ from 0° to 90°) ────────
  % x = +0.707·R·sinθ,  y = R·(−0.707·sinθ·sintilt + cosθ·costilt)
  \draw[EHAccretionGold, thick] plot[smooth, tension=0.55] coordinates {
    (0,                          {\R*\costilt})
    ({0.7071*\R*sin(15)}, {(-0.7071*sin(15)*\sintilt + cos(15)*\costilt)*\R})
    ({0.7071*\R*sin(30)}, {(-0.7071*sin(30)*\sintilt + cos(30)*\costilt)*\R})
    ({0.7071*\R*sin(45)}, {(-0.7071*sin(45)*\sintilt + cos(45)*\costilt)*\R})
    ({0.7071*\R*sin(60)}, {(-0.7071*sin(60)*\sintilt + cos(60)*\costilt)*\R})
    ({0.7071*\R*sin(75)}, {(-0.7071*sin(75)*\sintilt + cos(75)*\costilt)*\R})
    ({0.7071*\R},               {(-0.7071*\sintilt + 0*\costilt)*\R})
  };

  % ── Triangle edge: C→A (equator arc, φ from 315° to 225°) ─────────
  \draw[EHAccretionGold, thick]
    (\Cx, \Cy) arc (315:225:{\R} and {\eqb});

  % ── Parallel-transport vectors at vertex A ─────────────────────────
  %
  % Initial: pointing north = (0,0,1) in embedding.
  %   Screen: dx=0, dy = costilt·R.  Direction: straight up.
  %
  % After transport A→B→C→A, the vector arrives as the +φ direction
  % at A: embedding (0.707, −0.707, 0).
  %   Screen: dx = 0.707·R, dy = −0.707·sintilt·R.
  %
  \def\veclen{0.6}

  % Initial vector (blue): points north = up on screen
  \draw[->, EHJetBlue, very thick]
    (\Ax, \Ay) -- ++(0, \veclen)
    node[left, inner sep=2pt, font=\scriptsize, text=EHJetBlue] {initial};

  % Final vector (crimson): points along equator (+φ direction at A)
  % Direction (0.707, −0.707·sintilt) normalised, scaled to veclen
  \pgfmathsetmacro{\fdx}{0.7071}
  \pgfmathsetmacro{\fdy}{-0.7071*\sintilt}
  \pgfmathsetmacro{\flen}{sqrt(\fdx*\fdx + \fdy*\fdy)}
  \draw[->, EHHorizonCrimson, very thick]
    (\Ax, \Ay) -- ++({\veclen*\fdx/\flen}, {\veclen*\fdy/\flen})
    node[below, inner sep=2pt, font=\scriptsize, text=EHHorizonCrimson] {final};

  % ── Rotation arc showing π/2 ──────────────────────────────────────
  % Arc from final-vector direction to initial-vector direction
  % Final is at screen angle atan2(fdy, fdx); initial is at 90°.
  \pgfmathsetmacro{\arcrad}{0.35}
  \pgfmathsetmacro{\finalangle}{atan2(\fdy, \fdx)}
  \draw[EHHorizonCrimson, thin, ->]
    ({\Ax + \arcrad*\fdx/\flen}, {\Ay + \arcrad*\fdy/\flen})
    arc ({\finalangle}:90:{\arcrad})
    node[pos=0.45, right, inner sep=1pt, font=\scriptsize,
         text=EHHorizonCrimson] {$\pi/2$};

  % ── Vertex dots and labels ─────────────────────────────────────────
  \fill (\Ax, \Ay) circle (1.5pt);
  \fill (\Bx, \By) circle (1.5pt);
  \fill (\Cx, \Cy) circle (1.5pt);

  \node[font=\scriptsize, anchor=north east]
    at ({\Ax - 0.06}, {\Ay - 0.06}) {$A$};
  \node[font=\scriptsize, anchor=south]
    at (\Bx, {\By + 0.08}) {$B$};
  \node[font=\scriptsize, anchor=north west]
    at ({\Cx + 0.06}, {\Cy - 0.06}) {$C$};

  % ── Sphere label ───────────────────────────────────────────────────
  \node[font=\small] at (0, {-\R - 0.35}) {$S^2$};

\end{tikzpicture}

  \caption[Parallel transport around a closed loop on the two-sphere]{Parallel transport around a closed loop on the two-sphere
    demonstrates the path dependence that defines curvature.  A vector
    initially pointing north (blue) is parallel-transported around the
    spherical triangle $A \to B \to C \to A$, whose three sides are
    great-circle arcs (gold).  Upon returning to the starting vertex, the
    vector has rotated by $\pi/2$ to point along the equator (crimson) ---
    an angle equal to the solid angle enclosed by the triangle, here
    one-eighth of the sphere.  On a flat surface, the vector would return
    to its original orientation.  The rotation angle per unit area, in the
    limit of a vanishingly small loop, is the Gaussian curvature
    $K = 1/R^2$.}
  \label{fig:parallel-transport-loop}
\end{figure}

\paragraph{The Riemann tensor via the commutator.}

The covariant derivative $\cd{\mu}$ from \cref{sec:dg-covariant} --- the
torsion-free Levi-Civita connection --- was constructed so that
$\cd{\mu} V^\nu$ transforms as a tensor.  For partial
derivatives, the order of differentiation doesn't matter:
$\pd{\mu}\,\pd{\nu} f = \pd{\nu}\,\pd{\mu} f$ for any smooth function $f$.
For covariant derivatives acting on vectors, this commutativity fails on curved
manifolds.  The failure is measured by the \emph{Riemann curvature tensor},
defined through the commutator
%
\begin{equation}
  (\cd{\mu}\,\cd{\nu} - \cd{\nu}\,\cd{\mu})\, V^\rho
  = \riemann{\rho}{\sigma}{\mu}{\nu}\, V^\sigma \,.
  \label{eq:riemann-def}
\end{equation}
%
The left-hand side is the commutator $[\cd{\mu}, \cd{\nu}]$ applied to an
arbitrary vector field $V^\rho$.  The right-hand side defines a $(1,3)$
tensor $\riemann{\rho}{\sigma}{\mu}{\nu}$ that depends only on the connection,
not on the vector being transported.  If the Riemann tensor vanishes
everywhere, covariant derivatives commute, parallel transport is
path-independent, and the manifold is flat.  Conversely, a nonzero Riemann
tensor is the infinitesimal version of the rotation described in the opening
paragraph: it measures the curvature that a vanishingly small loop encloses.
\warnnote{The index ordering and sign convention for the Riemann tensor vary
across textbooks.  We follow Carroll \cite{Carroll2004} and
Wald \cite{Wald1984}: the upper index is first, and
\cref{eq:riemann-def} has no minus sign.  MTW \cite{MTW1973} uses the same
sign but different index ordering.}

\paragraph{Component formula.}

Expanding the commutator using \cref{eq:covariant-deriv-vector} and collecting
terms yields the component expression
%
\begin{equation}
  \riemann{\rho}{\sigma}{\mu}{\nu}
  = \pd{\mu}\,\christoffel{\rho}{\nu}{\sigma}
  - \pd{\nu}\,\christoffel{\rho}{\mu}{\sigma}
  + \christoffel{\rho}{\mu}{\lambda}\,\christoffel{\lambda}{\nu}{\sigma}
  - \christoffel{\rho}{\nu}{\lambda}\,\christoffel{\lambda}{\mu}{\sigma} \,.
  \label{eq:riemann-components}
\end{equation}
%
The first two terms involve derivatives of the Christoffel symbols --- and
since the Christoffel symbols are built from first derivatives of the metric
(\cref{eq:christoffel-formula}), the Riemann tensor involves \emph{second}
derivatives of the metric.  The last two terms are quadratic in the Christoffel
symbols: schematically, $R \sim \partial\Gamma + \Gamma\Gamma$.  Curvature is
a second-order property of the metric, invisible to the Christoffel symbols
alone.
\histnote{Riemann introduced this tensor in his 1854 Habilitationsschrift,
the same lecture where he defined manifolds.  Christoffel independently
rediscovered it in 1869 in the context of differential invariants.}

\paragraph{Symmetries and independent components.}

The Riemann tensor with all indices lowered,
$R_{\rho\sigma\mu\nu} = g_{\rho\alpha}\,\riemann{\alpha}{\sigma}{\mu}{\nu}$,
possesses three algebraic symmetries:
%
\begin{align}
  R_{\rho\sigma\mu\nu} &= -R_{\rho\sigma\nu\mu} \,,
  \label{eq:riemann-antisym1} \\
  R_{\rho\sigma\mu\nu} &= -R_{\sigma\rho\mu\nu} \,,
  \label{eq:riemann-antisym2} \\
  R_{\rho\sigma\mu\nu} &= R_{\mu\nu\rho\sigma} \,.
  \label{eq:riemann-pairsym}
\end{align}
%
The first two state that the tensor is antisymmetric in each index pair
separately: swapping either the first two or the last two indices flips the
sign.  The third states that the two pairs can be exchanged as a unit.  A
fourth relation, the \emph{algebraic Bianchi identity}, provides an
additional constraint:
%
\begin{equation}
  R_{\rho\sigma\mu\nu} + R_{\rho\mu\nu\sigma} + R_{\rho\nu\sigma\mu} = 0 \,.
  \label{eq:bianchi-algebraic}
\end{equation}
%
This is a cyclic sum over the last three indices.  Together, these four
symmetries reduce the number of independent components from the naive
$n^4 = 256$ (for $n = 4$) to
%
\begin{equation}
  \frac{n^2(n^2 - 1)}{12} = 20 \qquad (n = 4) \,.
  \label{eq:riemann-count}
\end{equation}
%
The count follows from treating the antisymmetric pairs $[\rho\sigma]$ and
$[\mu\nu]$ as symmetric-matrix indices: each pair has $n(n-1)/2 = 6$ values,
so the pair-symmetric tensor has $6 \times 7/2 = 21$ components before the
algebraic Bianchi identity removes one more.
\physnote{In two dimensions, \cref{eq:riemann-count} gives
$2^2(2^2 - 1)/12 = 1$: a two-dimensional manifold has exactly one independent
curvature component, the Gaussian curvature.  In three dimensions there are 6,
the same as the number of independent Ricci-tensor components, so the Ricci
tensor captures all the curvature.  Only in four and higher dimensions does
the Riemann tensor contain information beyond the Ricci tensor.}

\paragraph{The Ricci tensor and Ricci scalar.}

Contracting the Riemann tensor on its first and third indices defines the
\emph{Ricci tensor}:
%
\begin{equation}
  \ricci = \riemann{\lambda}{\mu}{\lambda}{\nu} \,.
  \label{eq:ricci-def}
\end{equation}
%
The Ricci tensor inherits the symmetry $R_{\mu\nu} = R_{\nu\mu}$ from the
pair symmetry of the Riemann tensor (\cref{eq:riemann-pairsym}), so it has 10
independent components in four dimensions --- the same count as the metric.
This is not a coincidence: the Einstein field equations equate a combination of
the Ricci tensor and metric to the stress-energy tensor, and both sides must
have the same number of independent components.

A further contraction with the inverse metric yields the \emph{Ricci scalar}
(or \emph{scalar curvature}):
%
\begin{equation}
  \ricciscalar = \inversemetric\, R_{\mu\nu} \,.
  \label{eq:ricci-scalar-def}
\end{equation}
%
This is a single scalar field that summarises the trace of the curvature at
each point.  For a maximally symmetric space (one with the most Killing vectors
possible), the Ricci scalar is constant and determines the geometry completely:
positive $\ricciscalar$ for spheres, zero for flat space, negative for
hyperbolic space.

\paragraph{The Einstein tensor and the Bianchi identity.}

The Riemann tensor satisfies a differential identity --- the
\emph{differential Bianchi identity}:
%
\begin{equation}
  \cd{\alpha}\, \riemann{\rho}{\sigma}{\mu}{\nu}
  + \cd{\mu}\, \riemann{\rho}{\sigma}{\nu}{\alpha}
  + \cd{\nu}\, \riemann{\rho}{\sigma}{\alpha}{\mu} = 0 \,.
  \label{eq:bianchi-differential}
\end{equation}
%
This is a cyclic sum of covariant derivatives over three indices.  Contracting
twice --- once on $\rho$ and $\mu$, then on $\sigma$ and $\alpha$ with the
inverse metric --- produces a remarkable result: a specific combination of the
Ricci tensor and scalar is automatically divergence-free.  This is the
\emph{contracted Bianchi identity}:
%
\begin{equation}
  \nabla^\mu\, \einstein = 0 \,,
  \label{eq:contracted-bianchi}
\end{equation}
%
where the \emph{Einstein tensor} is defined by
%
\begin{equation}
  \einstein = \ricci - \tfrac{1}{2}\,\metric\,\ricciscalar \,.
  \label{eq:einstein-tensor}
\end{equation}
%
The contracted Bianchi identity guarantees that the Einstein tensor is
divergence-free purely as a consequence of the geometry --- it holds for any
metric, regardless of any field equations.  This is the reason the Einstein
tensor appears on the left-hand side of Einstein's field equations:
setting $\einstein = 8\pi\,\stressenergy$ automatically
enforces $\nabla^\mu \stressenergy = 0$, the covariant statement of
energy-momentum conservation.
\xrefnote{The Einstein field equations and their physical content are the
subject of \cref{ch:general-relativity}.}

\paragraph{The Weyl tensor.}

The Riemann tensor can be decomposed into a trace part (determined by the Ricci
tensor) and a trace-free part called the \emph{Weyl tensor}
$\weyl{\rho}{\sigma}{\mu}{\nu}$.  In four dimensions, the decomposition is
%
\begin{equation}
  R_{\rho\sigma\mu\nu}
    = C_{\rho\sigma\mu\nu}
    + g_{\rho[\mu}\,R_{\nu]\sigma}
    - g_{\sigma[\mu}\,R_{\nu]\rho}
    - \tfrac{1}{3}\,\ricciscalar\;
      g_{\rho[\mu}\,g_{\nu]\sigma} \,,
  \label{eq:weyl-decomposition}
\end{equation}
%
where the square brackets denote antisymmetrisation:
$g_{\rho[\mu}\,R_{\nu]\sigma} = \tfrac{1}{2}(g_{\rho\mu}\,R_{\nu\sigma}
- g_{\rho\nu}\,R_{\mu\sigma})$.  The decomposition separates curvature into a
piece controlled by local matter (the Ricci terms, determined by the
stress-energy tensor through Einstein's equations) and a piece that propagates
freely (the Weyl tensor).  The Weyl tensor has 10 independent components in
four dimensions --- it vanishes identically for $n \leq 3$, where the Ricci
tensor already captures all curvature --- and shares all the symmetries of the
Riemann tensor plus the
additional property that every trace vanishes:
$C^\lambda{}_{\sigma\lambda\nu} = 0$.

In vacuum, where $R_{\mu\nu} = 0$ and $\ricciscalar = 0$, the Weyl tensor
\emph{is} the Riemann tensor.  It encodes the free gravitational field ---
the tidal distortions that propagate as gravitational waves and that produce
the image distortions computed by the ray tracer.
\physnote{The Weyl tensor is conformally invariant: multiplying the metric by a
smooth positive function leaves $\weyl{\rho}{\sigma}{\mu}{\nu}$ unchanged.
This makes it the natural curvature tensor for gravitational lensing, where the
conformal structure (light-cone shape) matters but the overall scale does not.}

\paragraph{A worked example: curvature of $S^2$.}

The two-sphere has been our running example throughout this chapter.  With
the Christoffel symbols from \cref{eq:sphere-christoffel} and the count from
\cref{eq:riemann-count} giving one independent curvature component in two
dimensions, a single computation determines the full Riemann tensor.
Substituting into \cref{eq:riemann-components}, the second term vanishes
because $\christoffel{\theta}{\theta}{\varphi} = 0$ and the third because
$\christoffel{\theta}{\theta}{\lambda} = 0$ for both values of $\lambda$;
only the first and fourth terms survive:
%
\begin{equation}
  \riemann{\theta}{\varphi}{\theta}{\varphi}
  = \pd{\theta}\,\christoffel{\theta}{\varphi}{\varphi}
    - \christoffel{\theta}{\varphi}{\varphi}\,
      \christoffel{\varphi}{\theta}{\varphi}
  = (\sin^2\!\theta - \cos^2\!\theta) + \cos^2\!\theta
  = \sin^2\!\theta \,,
  \label{eq:sphere-riemann}
\end{equation}
%
where the two intermediate results follow from
$\pd{\theta}(-\sin\theta\cos\theta) = \sin^2\!\theta - \cos^2\!\theta$
and $-(-\sin\theta\cos\theta)(\cot\theta) = \cos^2\!\theta$.

The Ricci tensor has components $R_{\theta\theta} = 1$ and
$R_{\varphi\varphi} = \sin^2\!\theta$; tracing with the inverse metric gives
the Ricci scalar
%
\begin{equation}
  \ricciscalar = g^{\theta\theta}\,R_{\theta\theta}
    + g^{\varphi\varphi}\,R_{\varphi\varphi}
  = \frac{1}{R^2} + \frac{1}{R^2} = \frac{2}{R^2} \,.
  \label{eq:sphere-ricci-scalar}
\end{equation}
%
\warnnote{The symbol $R$ is momentarily overloaded: $\ricciscalar$ on the
left is the Ricci scalar, while $R$ in the denominator is the sphere radius
from \cref{sec:dg-metric}.  Context and position in the equation resolve
the ambiguity.}
The factor of $R^2$ enters through the inverse metric, even though the
Christoffel symbols and the Riemann component
$\riemann{\theta}{\varphi}{\theta}{\varphi}$ are $R$-independent.  A larger
sphere is less curved: the Earth's surface, with $R \approx 6{,}400\;\text{km}$,
has $\ricciscalar \approx 5 \times 10^{-14}\;\text{m}^{-2}$ --- small enough
that a city appears flat, but detectable over continental distances.  The
Gaussian curvature $K = \ricciscalar/2 = 1/R^2$ recovers the standard result
from classical differential geometry.
\histnote{Gauss proved in his \emph{Theorema Egregium} (1827) that the
Gaussian curvature depends only on the intrinsic metric, not on how the surface
is embedded in three-dimensional space --- the first result in what would become
Riemannian geometry.}

The curvature tensors defined here complete the geometric toolkit of this
chapter.  The Riemann tensor encodes the full curvature information (20
independent components in four dimensions); the Ricci tensor (10 components)
captures how matter and energy curve spacetime; and the Weyl tensor (10
components) captures the remaining free-field curvature responsible for
gravitational lensing and tidal distortions.  With the Einstein tensor in
hand --- a divergence-free combination of the Ricci tensor and scalar that can
be consistently equated to a matter source --- the next chapter writes down
Einstein's field equations and begins to populate these geometric structures
with physics.
